%=====================================================================
\chapter{Markov Decision Processes}\label{ch:mdp}
%=====================================================================
\locator{5}{Part V --- Sequences, Decisions and Reinforcement}

\begin{outcomes}
By the end of this chapter you will be able to:
\begin{tight}
  \item \textbf{Formulate} a sequential decision problem as a Markov decision process:
        states, actions, transitions, rewards and discount.
  \item \textbf{Compute} the discounted return of a sequence of rewards, and
        \textbf{explain} the role of $\gamma$.
  \item \textbf{Define} state values and action values, and \textbf{write} the Bellman
        equations for them.
  \item \textbf{Perform} value iteration by hand for a small MDP and \textbf{read off}
        the optimal policy.
  \item \textbf{Describe} policy iteration and \textbf{compare} it with value iteration.
\end{tight}
\end{outcomes}

\begin{prereq}
\chref{hmm} (Markov chains and transition matrices) and, from \emph{Artificial
Intelligence}, agents and search. An MDP is a Markov chain in which an agent chooses
the transition probabilities by its actions, and is paid for doing so.
\end{prereq}

\begin{keyterms}
agent $\bullet$ environment $\bullet$ state $\bullet$ action $\bullet$ reward
$\bullet$ Markov decision process $\bullet$ transition model $\bullet$ policy
$\bullet$ return $\bullet$ discount factor $\gamma$ $\bullet$ state-value function
$V$ $\bullet$ action-value function $Q$ $\bullet$ Bellman expectation equation
$\bullet$ Bellman optimality equation $\bullet$ value iteration $\bullet$ policy
evaluation $\bullet$ policy improvement $\bullet$ policy iteration
\end{keyterms}

\section{Decisions Whose Consequences Unfold}

Classification makes one decision per example, and the next example does not
depend on it. Many project decisions are not like that. Refactoring ships nothing
this sprint but keeps the team fast in later ones; adding a developer to a late
project slows it before it speeds it up; releasing on a Friday saves a day now and
risks a weekend of firefighting. The value of an action depends on the situations
it leads to, and on the actions that will be taken there. The framework for such
problems is the \term{Markov decision process}, and it is the foundation of
reinforcement learning (\chref{rl}).

\dsfig{%
\begin{tikzpicture}[font=\scriptsize,
  box/.style={rounded corners=4pt, draw=#1, line width=1.1pt, fill=#1!10, text=#1!55!black,
              minimum width=2.8cm, minimum height=1.2cm, font=\small\bfseries},
  a/.style={-{Stealth[length=2.4mm]}, line width=1.1pt}]
\node[box=tealhl] (ag) at (0,0) {Agent};
\node[box=goldhl] (en) at (6.2,0) {Environment};
\draw[a, tealhl!70!black] (ag.north) .. controls +(0,1.3) and +(0,1.3) .. node[above, font=\scriptsize]
     {action $a_t$} (en.north);
\draw[a, goldhl!80!black] (en.south) .. controls +(0,-1.3) and +(0,-1.3) .. node[below, font=\scriptsize,
     align=center] {next state $s_{t+1}$, reward $r_{t+1}$} (ag.south);
\node[anchor=west, align=left, text width=5.2cm, font=\scriptsize, text=gray!55!black] at (8.2,0)
     {At each step the agent observes the state, chooses an action, and the environment
      responds with a new state and a numeric reward. The agent's goal is to choose
      actions that maximise the \emph{total} reward over time, not the next reward
      alone.};
\end{tikzpicture}}%
{The agent--environment loop. Every problem in Part V's last two chapters has this
shape.}{agent-env}

\section{The Markov Decision Process}

\begin{definitionbox}[Markov Decision Process]
An MDP consists of:
\begin{tight}
  \item a set of \term{states} $S$ and a set of \term{actions} $A$;
  \item a \term{transition model} $P(s' \mid s, a)$ --- the probability that action
        $a$ in state $s$ leads to state $s'$;
  \item a \term{reward} $R(s, a)$ received for taking $a$ in $s$ (or $R(s, a, s')$);
  \item a \term{discount factor} $0 \le \gamma < 1$.
\end{tight}
It has the Markov property: what happens next depends only on the current state and
action. A \term{policy} $\pi$ says what to do in each state, $a = \pi(s)$; solving the
MDP means finding the policy that maximises the expected discounted return.
\end{definitionbox}

\dsfig{%
\begin{tikzpicture}[font=\scriptsize,
  st/.style={circle, draw=#1, line width=1.1pt, fill=#1!12, text=#1!55!black,
             minimum size=1.35cm, font=\scriptsize\bfseries, align=center},
  op/.style={-{Stealth[length=2mm]}, secondary, line width=0.9pt},
  rp/.style={-{Stealth[length=2mm]}, goldhl!85!black, line width=0.9pt, dashed},
  lb/.style={font=\tiny, fill=white, inner sep=1pt}]
\node[st=greenhl!70!black] (H) at (0,0) {Healthy};
\node[st=accent] (D) at (4.4,0) {Debt-\\laden};
\node[st=gray!70!black] (F) at (8.8,0) {Frozen};
\draw[op] (H) to[out=120, in=60, looseness=6] node[lb, above] {ship: 0.8, $+10$} (H);
\draw[op] (H) to[bend left=18] node[lb, above] {ship: 0.2, $+10$} (D);
\draw[op] (D) to[out=120, in=60, looseness=6] node[lb, above] {ship: 0.5, $+5$} (D);
\draw[op] (D) to[bend left=18] node[lb, above] {ship: 0.5, $+5$} (F);
\draw[op] (F) to[out=120, in=60, looseness=6] node[lb, above] {ship: 1.0, $0$} (F);
\draw[rp] (D) to[bend left=18] node[lb, below] {refactor: 1.0, $-4$} (H);
\draw[rp] (F) to[bend left=38] node[lb, below] {refactor: 1.0, $-15$} (H);
\draw[rp] (H) to[out=-120, in=-60, looseness=6] node[lb, below] {refactor: $-2$} (H);
\node[anchor=west, font=\tiny, text=gray!55!black, align=left] at (10.0,0.4)
     {\textcolor{secondary}{\bfseries solid}: ship features\\[2pt]
      \textcolor{goldhl!85!black}{\bfseries dashed}: refactor\\[2pt]
      label: probability, reward\\[2pt]
      one step $=$ one sprint;\\
      discount $\gamma = 0.9$};
\end{tikzpicture}}%
{The technical-debt MDP of the worked examples. Shipping features earns value but lets
debt build up; refactoring costs a sprint's value and restores the codebase.}{debt-mdp}

\noindent The codebase is \emph{Healthy} (H), \emph{Debt-laden} (D) --- every change
takes longer than it should --- or \emph{Frozen} (F), so fragile that each change
breaks something else and the team does nothing but fight fires. Each sprint the
team either ships features, earning value (10 when healthy, 5 when debt-laden,
nothing when frozen) but risking more debt, or spends the sprint refactoring, which
costs value (2, 4, or 15 for a rewrite of a frozen codebase) and restores health.

\begin{definitionbox}[Return and Discounting]
The \term{return} from time $t$ is the discounted sum of future rewards:
\[
G_t = r_{t+1} + \gamma\,r_{t+2} + \gamma^2 r_{t+3} + \dots = \sum_{k=0}^{\infty}\gamma^k\,r_{t+k+1}
\]
With $\gamma < 1$ the sum is finite even for tasks that never end, and a reward $k$ steps
away is worth $\gamma^k$ of the same reward now. A small $\gamma$ makes the agent
short-sighted; $\gamma$ close to 1 makes it patient.
\end{definitionbox}

\begin{worked}[Two Returns]
With $\gamma = 0.9$:
\begin{tight}
  \item Rewards $10, 10, 10, 5, 5, 0, 0, \dots$ (keep shipping until the codebase
        freezes): $G = 10 + 9 + 8.1 + 5(0.729) + 5(0.656) = 27.1 + 3.645 + 3.281 =
        \mathbf{34.0}$.
  \item Rewards $10, 10, -4, 10, 10, \dots$ (refactor once, then keep shipping): the
        first two give $19$; the refactoring costs $4 \times 0.81 = 3.24$; and a
        healthy codebase thereafter earns, at the least, $10(0.729 + 0.656 + \dots) =
        10 \times 0.729/(1 - 0.9) = 72.9$ if it never gathers debt again. The
        refactoring pays for itself many times over.
\end{tight}
The second calculation needs to know what happens after the refactoring --- which is
exactly what the value functions below capture.
\end{worked}

\section{Value Functions and the Bellman Equations}

\begin{definitionbox}[State and Action Values]
Under a policy $\pi$, the \term{state value} $V^{\pi}(s)$ is the expected return
starting from $s$ and following $\pi$; the \term{action value} $Q^{\pi}(s, a)$ is the
expected return from taking $a$ in $s$ and following $\pi$ afterwards. They satisfy
the \term{Bellman expectation equation}:
\[
V^{\pi}(s) = R(s, \pi(s)) + \gamma\sum_{s'} P(s' \mid s, \pi(s))\,V^{\pi}(s')
\]
--- the value of a state is the immediate reward plus the discounted value of where
you go next. The best achievable values satisfy the \term{Bellman optimality
equation}:
\[
V^{*}(s) = \max_{a}\Big[R(s, a) + \gamma\sum_{s'} P(s' \mid s, a)\,V^{*}(s')\Big] = \max_a Q^{*}(s, a)
\]
and the optimal policy is to act greedily with respect to them:
$\pi^{*}(s) = \arg\max_a Q^{*}(s, a)$.
\end{definitionbox}

\begin{conceptbox}[Why a Recursive Definition Is Useful]
The Bellman equation turns a sum over infinitely many futures into a relation between
neighbouring states. For a known policy it is a set of linear equations --- one per
state --- that can be solved directly. For the optimal values it is not linear
(because of the max), but it can be solved by repeatedly applying it, which is value
iteration.
\end{conceptbox}

\section{Value Iteration}

\begin{definitionbox}[Value Iteration]
Start with $V(s) = 0$ for every state. Repeat until the values stop changing: for
every state, set
\[
V(s) \leftarrow \max_a\Big[R(s, a) + \gamma\sum_{s'} P(s' \mid s, a)\,V(s')\Big]
\]
Then read off the policy: in each state, the action achieving the maximum. After $k$
sweeps, $V$ is the best return achievable in $k$ steps; as $k$ grows it converges to
$V^{*}$.
\end{definitionbox}

\begin{worked}[Value Iteration on the Technical-Debt MDP]
$\gamma = 0.9$; start from $V = (0, 0, 0)$ for (H, D, F).

\smallskip
\textbf{Sweep 1.} With all values 0, only immediate rewards count: ship gives 10, 5, 0
and refactor $-2$, $-4$, $-15$. $V = (10,\ 5,\ 0)$; best action everywhere: ship.

\textbf{Sweep 2.}
\begin{tight}
  \item H: ship $10 + 0.9(0.8 \times 10 + 0.2 \times 5) = 10 + 8.1 = 18.1$; refactor
        $-2 + 0.9 \times 10 = 7$. $\;V(H) = 18.1$.
  \item D: ship $5 + 0.9(0.5 \times 5 + 0.5 \times 0) = 7.25$; refactor $-4 + 0.9 \times 10
        = 5$. $\;V(D) = 7.25$.
  \item F: ship $0 + 0.9 \times 0 = 0$; refactor $-15 + 9 = -6$. $\;V(F) = 0$.
\end{tight}
\textbf{Sweep 3.}
\begin{tight}
  \item H: ship $10 + 0.9(0.8 \times 18.1 + 0.2 \times 7.25) = 24.34$; refactor $-2 +
        16.29 = 14.29$. $\;V(H) = 24.34$.
  \item D: ship $5 + 0.9(0.5 \times 7.25 + 0) = 8.26$; refactor $-4 + 0.9 \times 18.1 =
        \mathbf{12.29}$. \textbf{Refactoring now wins.}
  \item F: ship $0$; refactor $-15 + 16.29 = \mathbf{1.29}$. \textbf{Refactoring wins.}
\end{tight}
\textbf{Converged} (after about 200 sweeps): $V^{*} = (78.64,\ 66.78,\ 55.78)$, and the
optimal policy is \textbf{ship when healthy, refactor when debt-laden or frozen}.

\smallskip
\textbf{Check} with the Bellman equation for that policy:
$V(H) = 10 + 0.9(0.8V(H) + 0.2V(D))$ and $V(D) = -4 + 0.9V(H)$ give
$V(H) = 9.28/0.118 = 78.64$. For comparison, ``always ship'' is worth only
$V(H) = 41.56$: the codebase eventually freezes and delivers nothing for ever.
\end{worked}

\section{Policy Iteration}

Value iteration improves values and only implicitly improves the policy. \term{Policy
iteration} alternates two explicit steps:
\begin{tightnum}
  \item \textbf{Policy evaluation:} compute $V^{\pi}$ for the current policy exactly,
        by solving the Bellman expectation equations --- one linear equation per
        state.
  \item \textbf{Policy improvement:} make the policy greedy with respect to $V^{\pi}$:
        $\pi(s) \leftarrow \arg\max_a\big[R(s, a) + \gamma\sum_{s'}P(s' \mid s, a)V^{\pi}(s')\big]$.
\end{tightnum}
Stop when the policy no longer changes. Each improvement is guaranteed not to make any
state worse, and there are finitely many policies, so it terminates --- usually in a
handful of iterations. On the technical-debt MDP it starts from ``always ship''
($V(H) = 41.56$), improves once to ``refactor when debt-laden or frozen'', evaluates
that ($V(H) = 78.64$), finds nothing to change, and stops: two evaluations against about
two hundred sweeps of value iteration, each evaluation costlier.

\rowcolors{2}{white}{lightbg}
\begin{longtable}{@{}L{3.0cm}L{6.0cm}L{6.0cm}@{}}
\toprule
\rowcolor{primary!15}
\textbf{} & \textbf{Value iteration} & \textbf{Policy iteration} \\
\midrule
\endhead
Each iteration & One cheap sweep: a max over actions per state & A full policy
evaluation (solve $|S|$ linear equations) plus an improvement \\
Iterations needed & Many; converges gradually & Few; stops exactly \\
Best when & Many states, cheap sweeps & Fewer states, or a good starting policy \\
Needs the model? & Yes: $P$ and $R$ & Yes: $P$ and $R$ \\
\bottomrule
\end{longtable}

\begin{alertbox}[Both Methods Need the Model]
Value and policy iteration are \emph{planning} algorithms: they compute the optimal
policy from a known $P(s' \mid s, a)$ and $R(s, a)$. In most real problems nobody
writes those down --- the probability that a refactoring really clears the debt, that
a new hire speeds the team up, or that a client accepts a release is unknown. \chref{rl} removes the model: the agent learns the
same value functions from experience, by trying actions and observing what happens.
\end{alertbox}

\begin{pitfall}
\begin{tight}
  \item \textbf{Maximising the next reward.} The point of an MDP is that the best
        action now may earn less now.
  \item \textbf{Choosing $\gamma$ carelessly.} A small $\gamma$ ignores long-term
        consequences; $\gamma = 1$ can make values infinite in tasks that never end.
  \item \textbf{A state that is not Markov.} If the right action depends on history
        the state does not record, add that history to the state.
  \item \textbf{Stopping value iteration early and trusting the values.} The policy
        often settles long before the values do; check both.
  \item \textbf{Rewarding the wrong thing.} An agent maximises exactly the reward it
        is given. Reward the outcome you want, not a proxy it can game.
\end{tight}
\end{pitfall}

%---------------------------------------------------------------------
\begin{lab}[Value Iteration and Policy Iteration]
\begin{lstlisting}[language=Python]
import numpy as np

# --- the technical-debt MDP: states H, D, F; actions ship (0), refactor (1)
S, A = ["H", "D", "F"], ["ship", "refactor"]
P = np.zeros((2, 3, 3))              # P[a, s, s']
P[0, 0] = [0.8, 0.2, 0.0]            # ship when healthy
P[0, 1] = [0.0, 0.5, 0.5]            # ship when debt-laden
P[0, 2] = [0.0, 0.0, 1.0]            # ship when frozen: stays frozen
P[1, :, 0] = 1.0                     # refactor: back to healthy
R = np.array([[10.0, 5.0, 0.0],      # R[a, s]: ship
              [-2.0, -4.0, -15.0]])  #          refactor
gamma = 0.9

# --- value iteration -------------------------------------------------------
V = np.zeros(3)
for it in range(1, 1000):
    Q = R + gamma * P @ V            # Q[a, s]
    V_new = Q.max(axis=0)
    if it <= 3:
        print(f"iteration {it}: V = {V_new.round(3)}  "
              f"policy = {[A[a] for a in Q.argmax(axis=0)]}")
    if np.abs(V_new - V).max() < 1e-8:
        break
    V = V_new
print(f"converged after {it} iterations: V = {V.round(2)}")
print("optimal policy:", dict(zip(S, [A[a] for a in Q.argmax(axis=0)])))

# --- policy iteration: evaluate exactly, then improve ----------------------
policy = np.zeros(3, dtype=int)                  # start: always ship
while True:
    Pp = P[policy, np.arange(3)]                 # 3 x 3 under the policy
    Rp = R[policy, np.arange(3)]
    V = np.linalg.solve(np.eye(3) - gamma * Pp, Rp)   # V = R + g P V
    print("evaluate", [A[a] for a in policy], "->", V.round(2))
    new = (R + gamma * P @ V).argmax(axis=0)     # improve greedily
    if (new == policy).all():
        break
    policy = new
\end{lstlisting}
Both methods reproduce the worked example. Now change \texttt{gamma} to 0.5 --- a
team that cares about this sprint far more than the ones after it --- and rerun. What
happens to the policy, and what does that say about teams that are always under
deadline pressure?
\end{lab}

%---------------------------------------------------------------------
\begin{checkpoint}
\begin{tightnum}
  \item With $\gamma = 0.5$, what is the return of the rewards $4, 4, 4, 0, 0, \dots$?
  \item What is the difference between $V^{\pi}(s)$ and $Q^{\pi}(s, a)$?
  \item In the technical-debt MDP, why does ``always ship'' score so much less than
        the optimal policy, when shipping always earns more immediately than
        refactoring?
\end{tightnum}
\end{checkpoint}

%---------------------------------------------------------------------
\begin{chaptersummary}
\begin{tight}
  \item An MDP has states, actions, a transition model, rewards and a discount; a
        policy maps states to actions.
  \item The return is the discounted sum of future rewards; $\gamma$ sets how far
        ahead the agent looks.
  \item The Bellman equations relate each state's value to its neighbours'; the
        optimal policy acts greedily with respect to the optimal action values.
  \item Value iteration repeatedly applies the Bellman optimality update; policy
        iteration alternates exact evaluation and greedy improvement and usually
        needs very few rounds.
  \item Both are planning methods that require the model; reinforcement learning
        learns without it.
\end{tight}
\end{chaptersummary}

\begin{reviewq}
  \item \textbf{(MCQ)} The discount factor $\gamma$ controls: (a)~the learning rate
        (b)~how much future rewards count relative to immediate ones (c)~the
        probability of a transition (d)~the number of actions.
  \item \textbf{(MCQ)} Value iteration requires: (a)~no model (b)~the transition and
        reward model (c)~labelled examples (d)~a neural network.
  \item \textbf{(MCQ)} The optimal policy chooses in each state the action that
        maximises: (a)~$R(s, a)$ (b)~$Q^{*}(s, a)$ (c)~$V^{*}(s)$ (d)~$P(s' \mid s, a)$.
  \item \textbf{(Short)} Explain the Bellman expectation equation in words.
  \item \textbf{(Short)} Why does policy iteration usually need fewer iterations than
        value iteration?
  \item \textbf{(Short)} Give an example of a reward that an agent could satisfy in an
        unintended way.
  \item \textbf{(Applied)} Perform sweep 4 of value iteration on the technical-debt
        MDP, starting from the sweep-3 values $(24.34, 12.29, 1.29)$.
  \item \textbf{(Applied)} In the technical-debt MDP, raise the cost of refactoring
        a debt-laden codebase to $-20$. Write the Bellman equations for the policy
        ``ship when healthy or debt-laden, refactor when frozen'' and solve them for
        $V(H)$.
\end{reviewq}
